\documentclass[11pt,a4paper]{article}

\usepackage[margin=1in]{geometry}
\usepackage{amsmath}
\usepackage{amssymb}
\usepackage{tabularx}
\usepackage{enumitem}
\usepackage[hidelinks]{hyperref}

\setlist[itemize]{itemsep=1pt,topsep=3pt,parsep=0pt}
\setlist[enumerate]{itemsep=1pt,topsep=3pt,parsep=0pt}
\setlength{\parskip}{4pt}
\setlength{\parindent}{0pt}
\newcolumntype{Y}{>{\raggedright\arraybackslash}X}

\title{Advanced Natural Language Processing\\ Comprehensive Final Examination Study Notes}
\author{Course Notes}
\date{}

\begin{document}

\maketitle

\section{Part I: Recurrent Models}

\subsection{The Simple Recurrent Neural Network (RNN)}

\paragraph{What it is.} A neural network with a loop: it processes a sequence one token at
a time and carries a hidden state vector from each step into the next, so that the state
acts as a summary of everything seen so far.

\paragraph{The idea, plainly.} Imagine reading a sentence with a single sticky note in your
hand. At each word you (i) read the word, (ii) read your sticky note, (iii) write a new
sticky note that combines the two, and (iv) throw the old one away. The sticky note is the
hidden state $h^{(t)}$. Because you rewrite it in exactly the same way at every word, you
can read a sentence of any length with a fixed amount of machinery.

\paragraph{Mechanics.} With input embedding $x^{(t)}$ and hidden state $h^{(t)}$:
\[
h^{(t)} = \tanh\!\left(U h^{(t-1)} + W x^{(t)}\right), \qquad
o^{(t)} = \mathrm{softmax}\!\left(V h^{(t)}\right).
\]
Three observations that the exam likes:
\begin{itemize}
\item $x$ and $h$ change at every time step; the parameters $U$, $W$, $V$ do \textbf{not}.
  They are shared across all positions.
\item $U$ mixes in history, $W$ mixes in the present word, $V$ reads out a prediction.
\item $h^{(0)}$ is usually the zero vector.
\end{itemize}

\paragraph{Problem it solves.} Variable-length input and order-sensitivity, using a constant
number of parameters. Weight sharing means a model trained on 10-word sentences can process
a 40-word sentence.

\paragraph{Shapes of RNN problems.} The same cell supports several task shapes
(Karpathy's taxonomy):
\begin{center}
\begin{tabularx}{\textwidth}{lY}
\textbf{Shape} & \textbf{Example} \\
\hline
one-to-one & Ordinary classification (no recurrence really needed). \\
one-to-many & Image captioning: one image in, a sentence out. \\
many-to-one & Sentiment analysis: the last hidden state represents the whole sentence. \\
many-to-many (aligned) & POS tagging / NER: one label per token. \\
many-to-many (unaligned) & Machine translation: read all, then write all. \\
\end{tabularx}
\end{center}

\paragraph{Advantages.}
\begin{itemize}
\item Handles arbitrary length with fixed parameter count.
\item Parameter sharing gives strong generalisation across positions.
\item Small memory footprint; state is one vector.
\item Naturally streaming/online --- it can process a token the moment it arrives.
\end{itemize}

\paragraph{Disadvantages.}
\begin{itemize}
\item \textbf{Sequential by construction}, so it cannot be parallelised along time. Step $t$
  needs $h^{(t-1)}$.
\item \textbf{Vanishing (and exploding) gradients} during backpropagation through time.
\item \textbf{It can only see the past}, not the future context of a token.
\item \textbf{It must forget the same amount of history at every step}, because the update
  is a single fixed $\tanh$ transformation --- there is no mechanism to say ``this token
  matters, keep it for 50 steps''.
\end{itemize}

\paragraph{What it cannot solve.} Long-term dependencies. Theoretically a simple RNN is
capable of representing them --- you could hand-pick weights that do it --- but in practice
gradient descent does not find those weights. This gap between representational capacity
and learnability is a favourite exam distinction.

\paragraph{What it gave rise to.} The vanishing gradient problem, which motivated gating
(LSTM, GRU); the past-only limitation, which motivated bidirectional RNNs; the sequential
bottleneck, which eventually motivated the Transformer.

\paragraph{Use it when.} Short sequences, streaming/low-latency settings, very small
compute or memory budgets, or as a cheap baseline.

\paragraph{Do not use it when.} Dependencies span more than roughly 10--20 tokens, you have
GPUs and want to exploit them, or you need state-of-the-art accuracy.

\paragraph{Counterfactual.} \emph{What if the weights were not shared across time steps?}
You would need a separate $U,W,V$ per position. The parameter count would grow with sequence
length, the model could never process a sequence longer than the longest one it was built
for, and a pattern learned at position 3 would not transfer to position 30. Weight sharing
is not an optimisation detail --- it is the entire reason an RNN generalises over sequences.

\subsection{Backpropagation Through Time and the Vanishing Gradient}

\paragraph{What it is.} BPTT is ordinary backpropagation applied to the RNN after
``unrolling'' it: the loop is expanded into a deep feed-forward network with one layer per
time step, all layers sharing the same weights.

\paragraph{The idea, plainly.} A 50-word sentence turns the RNN into a 50-layer network.
Gradients must be carried back through all 50 layers by repeated multiplication.

\paragraph{Mechanics, and why it fails.} The gradient of the loss at step $T$ with respect
to the state at step $t$ contains a product of $T-t$ Jacobians:
\[
\frac{\partial h^{(T)}}{\partial h^{(t)}} = \prod_{k=t+1}^{T} \frac{\partial h^{(k)}}{\partial h^{(k-1)}}.
\]
Each factor involves $U$ and the derivative of $\tanh$ (which is at most 1, and usually
much less). Multiplying many numbers smaller than one gives a compounding shrinkage:
\begin{itemize}
\item If the factors are $<1$, the gradient decays roughly geometrically. After 30--50
  steps it is numerically zero: the early time steps receive no learning signal.
  This is the \textbf{vanishing gradient problem}.
\item If the factors are $>1$, the gradient blows up: the \textbf{exploding gradient
  problem}. This one is easy to patch with gradient clipping. Vanishing is the hard one,
  because you cannot ``unclip'' a zero.
\end{itemize}

\paragraph{Why this is a modelling problem, not a numerical nuisance.} A zero gradient at
early time steps means the model is literally unable to learn that the word ``France''
25 tokens ago should influence the current prediction. The failure is in learning long-term
dependencies, exactly the property we needed.

\paragraph{Counterfactual.} \emph{What if the state update were additive instead of
multiplicative?} Then the Jacobian would be close to the identity and the gradient would
flow back nearly unchanged. This is precisely the design that LSTMs adopt for the cell
state, and later the design that residual connections adopt in Transformers. The same trick
appears twice in this syllabus --- notice it.

\subsection{Bidirectional RNNs}

\paragraph{What it is.} Two independent RNNs over the same sequence, one running left to
right and one right to left, whose hidden states are concatenated at each position.

\paragraph{The idea, plainly.} To tag the word ``bank'' you would like to see both ``river''
before it and ``flooded'' after it. One RNN reads the sentence forwards, another reads it
backwards, and each position ends up with a representation of its full context:
$h^{(t)} = [\overrightarrow{h}^{(t)} ; \overleftarrow{h}^{(t)}]$.

\paragraph{Problem it solves.} The ``can only see the past'' limitation of the simple RNN.

\paragraph{Advantages.} Substantially better on tagging tasks (POS, NER, chunking) because
disambiguation usually needs right context. A bidirectional GRU is a strong default
starting point for sequence tagging.

\paragraph{Disadvantages.} Roughly doubles compute and parameters. Crucially, it needs the
\emph{entire} sequence before it can produce any output.

\paragraph{What it cannot solve.} It does not help with generation or with streaming. It
does not fix vanishing gradients --- each direction still suffers from them.

\paragraph{Use it when.} Offline tasks where the full input is available: classification,
tagging, encoding the source side of a translation system.

\paragraph{Do not use it when.} Real-time/streaming decoding, or any autoregressive
generation task --- the backward pass would let the model see the future tokens it is
supposed to predict, which is leakage.

\paragraph{Counterfactual.} \emph{What if we made a language model bidirectional?}
It would be trivially able to cheat: predicting token $t$ while its backward pass has
already read token $t$. This is exactly why BERT cannot simply be trained left-to-right
bidirectionally and instead has to invent \emph{masking} (Section~\ref{sec:bert}). The
bidirectional/generation tension is a recurring exam theme.

\section{Part II: Encoder--Decoder and Attention}

\subsection{The Sequence-to-Sequence (Encoder--Decoder) Model}

\paragraph{What it is.} Two recurrent networks joined together. An \textbf{encoder} RNN
reads the entire input sequence and compresses it into a fixed-size vector called the
\emph{context vector} (or thought vector). A \textbf{decoder} RNN is initialised with that
vector and generates the output sequence one token at a time.

\paragraph{The idea, plainly.} Translation is not a tagging problem: ``I am hungry'' (3
words) becomes ``J'ai faim'' (2 words), and the word order need not match. So you cannot
emit one output per input. Instead: read the whole English sentence, form a single mental
summary, then write the French sentence from that summary. The encoder is the reader, the
context vector is the summary, the decoder is the writer.

\paragraph{Mechanics.}
\begin{enumerate}
\item Encoder consumes $x_1 \ldots x_n$; keep the final hidden state $c = h^{(n)}_{\text{enc}}$.
\item Initialise decoder state with $c$; feed a start token \texttt{<sos>}.
\item At each decoder step, produce a distribution over the vocabulary, emit a token, and
  feed that token back in as the next input.
\item Stop at \texttt{<eos>}.
\end{enumerate}
During training the decoder is fed the \emph{ground-truth} previous token rather than its
own prediction (\textbf{teacher forcing}); at inference it must consume its own outputs
(\textbf{autoregressive generation}). The mismatch between these two regimes is called
\emph{exposure bias}.

\paragraph{Problem it solves.} Sequence transduction with unequal, unaligned input and
output lengths: machine translation, summarisation, dialogue, speech recognition,
question answering.

\paragraph{Advantages.} Fully general over length pairs; end-to-end trainable; no need for
hand-built alignment or phrase tables as in statistical MT.

\paragraph{Disadvantages.}
\begin{itemize}
\item \textbf{The information bottleneck.} The whole source sentence, however long, must fit
  into one fixed-size vector. Quality degrades sharply as sentences get longer.
\item The encoder's final state is biased toward the end of the source sentence.
\item Still sequential on both sides; still suffers exposure bias.
\item Errors compound during generation: one wrong token corrupts everything after it.
\end{itemize}

\paragraph{What it cannot solve.} Long inputs. A 50-word legal sentence and a 5-word
greeting get the same number of bits of context.

\paragraph{What it gave rise to.} Attention, invented precisely to remove this bottleneck.

\paragraph{Use it when.} Any input-sequence to output-sequence mapping where the two are not
aligned position-by-position.

\paragraph{Do not use it when.} The task is aligned tagging (use a bidirectional encoder
with per-token outputs; a decoder is wasted machinery) or classification (use the encoder
alone).

\paragraph{Counterfactual.} \emph{What if the context vector were passed to every decoder
step instead of only initialising the first one?} This is in fact a common variant and it
helps a little --- the decoder stops ``forgetting'' the source as it generates. But it does
not fix the underlying problem, because the vector is still a single fixed-size summary.
The real fix has to let the decoder look back at \emph{all} encoder states, which is
attention.

\subsection{Attention}

\paragraph{What it is.} A mechanism that lets the decoder, at each output step, look back
over \emph{all} encoder hidden states and build a fresh, weighted summary tailored to the
word it is about to generate.

\paragraph{The idea, plainly.} A human translator does not memorise the source sentence and
then recite a translation. They glance back at the relevant source word for each word they
write. Attention gives the model the same ability: while producing the French word for
``hungry'' it puts most of its weight on the English token ``hungry'' and little elsewhere.

\paragraph{Mechanics.} Keep every encoder state $h_1 \ldots h_n$ (do not throw them away).
At decoder step $t$ with decoder state $s_{t}$:
\begin{align*}
e_{t,i} &= \mathrm{score}(s_{t}, h_i) && \text{relevance of source position } i \\
\alpha_{t,i} &= \frac{\exp(e_{t,i})}{\sum_j \exp(e_{t,j})} && \text{softmax: weights sum to 1} \\
c_t &= \sum_{i=1}^{n} \alpha_{t,i}\, h_i && \text{context vector for this step}
\end{align*}
The score function may be a dot product ($s^\top h$), a scaled dot product, or a small
feed-forward network (additive/Bahdanau attention). The context $c_t$ is then concatenated
with $s_t$ to predict the output token.

\paragraph{Problem it solves.} The fixed-size bottleneck. The context is now
\emph{dynamic}: a different summary for every output step, and its size effectively grows
with the input because all $n$ states remain available.

\paragraph{Advantages.}
\begin{itemize}
\item Large quality gains, especially on long sentences.
\item Creates a shortcut for gradients: the loss at output step $t$ connects directly to
  encoder state $h_i$ without travelling through $n$ recurrent steps. This substantially
  eases the vanishing gradient problem.
\item The weights $\alpha_{t,i}$ form a soft alignment matrix, which is interpretable
  (a partial answer to the black-box complaint).
\end{itemize}

\paragraph{Disadvantages.} Computes $n$ scores per output step, so cost is $O(n \cdot m)$
for input length $n$ and output length $m$; memory grows because all encoder states are
retained. Interpretability of attention weights is weaker than it looks --- attention is
correlation, not explanation.

\paragraph{What it cannot solve.} Attention on top of an RNN still leaves the RNN's
sequential computation intact. Training is still slow, because $h_i$ still has to be
computed one step at a time. This residual complaint is exactly what the Transformer
answers.

\paragraph{Use it when.} Any encoder--decoder task; effectively always.

\paragraph{Counterfactual.} \emph{What if we replaced the softmax with a hard argmax
(``hard attention'')?} The model would attend to exactly one source position. This is
cheaper at inference and sometimes more interpretable, but the operation is not
differentiable, so you would need reinforcement learning or a relaxation to train it.
Soft attention is used almost universally because it is differentiable end-to-end --- a
good illustration of a design chosen for trainability rather than realism.

\subsection{Why GPUs Matter (and Why Recurrence Became Unacceptable)}

\paragraph{The hardware argument.} CPUs are latency-optimised: a few very fast cores,
designed to finish a single instruction chain quickly. GPUs are bandwidth-optimised:
thousands of simple cores plus very fast VRAM, designed to perform the same operation on
enormous amounts of data at once.

\paragraph{Why deep networks fit GPUs.} A neural layer is a matrix multiplication --- many
identical neurons performing the same computation on different data. That is precisely the
workload a GPU is built for. The more arithmetic per unit of data movement, the larger the
GPU's advantage. Massive thread parallelism also hides memory latency.

\paragraph{Why this doomed the RNN.} An RNN's time axis is a strict dependency chain: you
cannot compute $h^{(5)}$ before $h^{(4)}$. So the one axis with the most work in it cannot
be parallelised. As GPUs got faster, RNN training time barely improved. A model that can
compute all positions simultaneously would convert hardware progress directly into research
progress. That model is the Transformer, and the observation above is the real motivation
behind ``Attention Is All You Need''.

\section{Part III: The Transformer}

\subsection{Overview}

\paragraph{What it is.} An encoder--decoder architecture that removes recurrence and
convolution entirely, and relies only on attention plus position-wise feed-forward layers.
Introduced in \emph{Attention Is All You Need} (Vaswani et al., 2017).

\paragraph{The idea, plainly.} If attention already lets a decoder look anywhere in the
source, why not let \emph{every} token look at every other token directly, and drop the
recurrent chain altogether? Then all positions can be computed in one parallel matrix
multiplication. The price is that, without recurrence, the model has no idea what order the
tokens came in --- so order must be injected separately.

\paragraph{The stack, in the original paper.}
\begin{itemize}
\item Encoder: $N = 6$ identical layers. Each layer = multi-head self-attention $\rightarrow$
  Add \& Norm $\rightarrow$ position-wise feed-forward network $\rightarrow$ Add \& Norm.
\item Decoder: $N = 6$ identical layers. Each layer = \emph{masked} multi-head self-attention
  $\rightarrow$ Add \& Norm $\rightarrow$ encoder--decoder (cross) attention $\rightarrow$
  Add \& Norm $\rightarrow$ feed-forward $\rightarrow$ Add \& Norm.
\item Dimensions: $d_{\text{model}} = 512$, $h = 8$ heads, $d_k = d_v = 64$,
  $d_{\text{ff}} = 2048$. Every sub-layer and the embedding layer outputs
  $d_{\text{model}}$, which is what makes residual connections possible.
\item Every sub-layer is wrapped as $\mathrm{LayerNorm}(x + \mathrm{Sublayer}(x))$.
\end{itemize}

\subsection{Self-Attention}

\paragraph{What it is.} Attention where the queries, keys and values all come from the same
sequence: each token attends to every token in its own sentence, including itself.

\paragraph{Why ``self''.} Ordinary attention asks ``which source words matter for this
target word?'' --- the query comes from a different sequence and often a different task.
Self-attention asks ``which words in \emph{this same} sentence matter for understanding
\emph{this} word?''. In ``The animal didn't cross the street because \textbf{it} was too
tired'', self-attention is what lets ``it'' bind to ``animal'' rather than ``street''.
Trivial for a human, historically very hard for a machine.

\paragraph{The Q/K/V metaphor (from information retrieval).} Each token is linearly
projected three times:
\begin{itemize}
\item \textbf{Query} --- the search string you type into YouTube.
\item \textbf{Key} --- the video titles the system matches your string against.
\item \textbf{Value} --- the actual videos returned to you.
\end{itemize}
Ranking by query--key similarity and returning a blend of the corresponding values is
exactly what we want: score every other word (key) against the current word (query), then
mix their contents (values) accordingly.

\paragraph{Mechanics.} With $X \in \mathbb{R}^{n \times d_{\text{model}}}$:
\[
Q = XW^Q,\quad K = XW^K,\quad V = XW^V, \qquad
\mathrm{Attention}(Q,K,V) = \mathrm{softmax}\!\left(\frac{QK^\top}{\sqrt{d_k}}\right)V .
\]
Step by step:
\begin{enumerate}
\item $QK^\top$ gives an $n \times n$ matrix of raw similarity scores (dot product = how
  aligned two vectors are).
\item Divide by $\sqrt{d_k}$ --- \textbf{scaling}. With $d_k = 64$ this is a division by 8.
\item Softmax each row so the weights over the sentence sum to 1.
\item Multiply by $V$ to get, for each token, a weighted blend of all tokens' values.
\end{enumerate}
Note that $d_k = 64$ is deliberately smaller than $d_{\text{model}} = 512$: each head works
in a compressed subspace, so that $h=8$ heads together cost about the same as one full-width
attention.

\paragraph{Problem it solves.} Direct, constant-length paths between any two positions. In
an RNN, information from token 1 reaches token 50 through 49 sequential transformations; in
self-attention it is one hop. It also makes the whole sequence computable in parallel,
because the operation is just three matrix multiplications and a softmax.

\paragraph{Advantages.}
\begin{itemize}
\item Maximum path length between any two tokens is $O(1)$: excellent for long-range
  dependencies.
\item Fully parallel across positions: converts GPU throughput directly into training speed.
\item Content-based, dynamic routing: what attends to what depends on the actual input, not
  on fixed positions as in a convolution.
\end{itemize}

\paragraph{Disadvantages.}
\begin{itemize}
\item $O(n^2)$ time and memory in the sequence length, because of the $n \times n$ score
  matrix. This is the single biggest cost driver in modern LLMs.
\item Permutation invariant on its own --- it has no notion of order.
\item No built-in locality bias, so it is data-hungry: it must learn from scratch that
  nearby words are usually more relevant.
\end{itemize}

\paragraph{Counterfactual 1.} \emph{What if we dropped the $\sqrt{d_k}$ scaling?} For large
$d_k$ the dot products have large variance and grow in magnitude. Feeding large values into
softmax pushes it into a saturated regime --- one weight near 1, the rest near 0 --- where
its gradient is almost zero. Training would stall or become unstable. The scaling keeps the
logits in a range where softmax has usable gradients. Expect this as an exam question.

\paragraph{Counterfactual 2.} \emph{What if $Q$, $K$ and $V$ were the same vectors, with no
projections?} The score matrix would be $XX^\top$, which is symmetric, so ``A attends to B''
would always equal ``B attends to A''. Real linguistic relations are asymmetric --- a
pronoun depends on its antecedent, not the reverse. Separate $W^Q$ and $W^K$ break the
symmetry. A separate $W^V$ additionally lets the model decide that the information worth
\emph{passing on} differs from the information used for \emph{matching}.

\subsection{Multi-Head Attention}

\paragraph{What it is.} Run $h$ self-attention operations in parallel with different learned
projections, concatenate their outputs, and pass the result through a final linear layer
$W^O$.
\[
\mathrm{MultiHead}(Q,K,V) = \mathrm{Concat}(\mathrm{head}_1, \ldots, \mathrm{head}_h)\,W^O,
\]
\[
\mathrm{head}_i = \mathrm{Attention}(XW_i^Q, XW_i^K, XW_i^V).
\]

\paragraph{The idea, plainly.} A single attention distribution has to commit to one
weighting. But ``it'' may need to attend to its antecedent \emph{and} the current verb
\emph{and} the sentence topic at the same time. One head cannot express three different
weightings simultaneously; averaging them muddles all three. Multiple heads let the model
maintain several relationship types in parallel --- one head may track syntactic subjects,
another coreference, another local adjacency. As the lecture puts it: multiple heads can
(hopefully) focus on more things than one head can.

\paragraph{Advantages.} Richer representations at essentially no extra cost, since each head
is narrower ($d_k = d_{\text{model}} / h$). Also gives redundancy: heads can be pruned later.

\paragraph{Disadvantages.} Many heads are demonstrably redundant and can be removed with
little loss --- so the mechanism is not as efficient as it appears. Head interpretations are
suggestive, not reliable.

\paragraph{Counterfactual.} \emph{What if we used one head of full width $d_{\text{model}}$
instead of eight of width 64?} Parameter count would be similar, but the layer could
represent only one attention pattern per layer. Empirically this is clearly worse. The
lesson is that the benefit comes from having \emph{multiple simultaneous} attention
distributions, not from extra parameters.

\subsection{Positional Encoding}

\paragraph{What it is.} A vector added to each input embedding that encodes the token's
position in the sequence.

\paragraph{Why it is necessary.} Self-attention treats its input as a \emph{set}. Shuffle
the tokens and every attention score is unchanged; the outputs just permute along with the
inputs. Without positional information the Transformer is a sophisticated bag of words, and
``dog bites man'' equals ``man bites dog''.

\paragraph{Mechanics (sinusoidal).} For position $k$, embedding dimension $d$, index
$0 \le i < d/2$, and a scalar $n$ set to $10000$ in the paper:
\[
PE(k, 2i) = \sin\!\left(\frac{k}{n^{2i/d}}\right), \qquad
PE(k, 2i+1) = \cos\!\left(\frac{k}{n^{2i/d}}\right).
\]
Properties worth memorising:
\begin{itemize}
\item Same dimension as the embedding, so it can simply be added.
\item Each dimension is a sinusoid of a different wavelength, from short (fast-varying, fine
  position) to very long (slow-varying, coarse position) --- like the hands of a clock at
  different speeds.
\item Because it is a fixed periodic function, the encoding of position 7 is the same in
  every example regardless of sentence length, and the model can extrapolate somewhat to
  lengths not seen in training.
\item $PE_{k+\delta}$ is a fixed linear function of $PE_k$, so relative offsets are easy for
  the model to compute.
\item Positional encodings can also be \emph{learned}; in practice learned embeddings did
  not perform better in the original experiments, and they cannot extrapolate beyond the
  trained maximum length.
\end{itemize}

\paragraph{Counterfactual 1.} \emph{What if we omitted positional encoding entirely?} The
model becomes permutation-invariant. It could still do bag-of-words tasks like topic
classification reasonably, but syntax, negation scope, and translation would collapse.

\paragraph{Counterfactual 2.} \emph{What if we used a simple integer position, $p_k = k$?}
The magnitude would grow without bound, dominating the embedding and destabilising training,
and the network would have to learn to decode a scalar into useful relative-distance
features. The sinusoidal scheme is bounded, smooth, and already exposes multiple scales.

\paragraph{Counterfactual 3.} \emph{What if position were concatenated instead of added?}
This works and is arguably cleaner conceptually, but it costs extra dimensions in every
subsequent layer. Addition is chosen for efficiency; the model can learn to allocate
subspaces to content versus position.

\subsection{Residual Connections and Layer Normalisation (Add \& Norm)}

\paragraph{What it is.} Around every sub-layer, the input is added to the output and the sum
is normalised: $\mathrm{LayerNorm}(x + \mathrm{Sublayer}(x))$.

\paragraph{The ``Add'' part.} The residual connection provides an identity path from input to
output. Two benefits:
\begin{itemize}
\item \textbf{Preserves information.} Attention re-mixes tokens; without the shortcut, the
  original token identity would be progressively washed out across 6 (or 96) layers.
\item \textbf{Handles vanishing gradients.} The derivative through $x + f(x)$ is
  $1 + f'(x)$, so gradient can flow straight back through the $1$. This is the same additive
  highway idea as the LSTM cell state --- and it is why 100-layer Transformers are trainable
  at all.
\end{itemize}
Concretely, in the first encoder layer the ``Add'' literally adds the multi-head attention
output to the (embedding + positional encoding) input.

\paragraph{The ``Norm'' part: layer normalisation.} Normalise each vector to zero mean and
unit variance \emph{across its own feature dimensions}, then rescale with learned gain and
bias. Contrast with batch normalisation, which computes mean and variance for each neuron
across a mini-batch of training cases. Batch norm dramatically reduces training time in
feed-forward networks, but it is a poor fit here:
\begin{itemize}
\item Sequences have variable length and heavy padding, so per-neuron batch statistics are
  noisy and ill-defined.
\item Batch norm behaves differently at training and inference time, and depends on batch
  size --- awkward for autoregressive decoding one token at a time.
\item Layer norm is computed per token, independently of the batch, so it behaves
  identically in training and inference.
\end{itemize}

\paragraph{Pre-norm vs post-norm.} The original paper normalises \emph{after} the residual
addition (post-norm). Modern large models (including LLaMA) normalise the \emph{input} of
each sub-layer instead (pre-norm), because post-norm requires careful learning-rate warm-up
and becomes unstable at great depth. Pre-norm keeps a completely clean residual path from
the first layer to the last.

\paragraph{Counterfactual.} \emph{What if we removed the residual connections?} Deep stacks
would become untrainable: gradients would vanish and the token's own identity would be lost
after a few rounds of attention mixing. Empirically, removing residuals from a Transformer
is catastrophic, far more so than removing multi-head structure.

\subsection{The Position-wise Feed-Forward Network}

\paragraph{What it is.} A two-layer MLP applied independently and identically to every
position:
\[
\mathrm{FFN}(x) = \max(0,\, xW_1 + b_1)W_2 + b_2 ,
\]
with inner dimension $d_{\text{ff}} = 2048$ against $d_{\text{model}} = 512$.

\paragraph{Why it is there.} Attention is essentially a weighted \emph{average} --- a linear
operation once the weights are fixed. Stacking averages without nonlinearity would collapse.
The FFN provides the per-token nonlinear processing: attention gathers information from
other tokens, the FFN then thinks about it. A useful slogan: \emph{attention mixes across
tokens, the FFN mixes across features}. It also holds most of the parameters (roughly
two-thirds of a Transformer block), and is widely believed to store much of the model's
factual knowledge --- which is why Mixture-of-Experts models make \emph{this} layer the
expert layer.

\subsection{The Decoder}

\paragraph{Three distinguishing features.}
\begin{enumerate}
\item \textbf{Right-shifted input.} The target sequence is offset by one position, so that
  when predicting token $i$ the decoder has been given tokens $1 \ldots i-1$. This is what
  makes autoregression well-defined.
\item \textbf{Masked multi-head self-attention.} Positions are forbidden from attending to
  \emph{subsequent} positions. Implemented by setting the corresponding scores to $-\infty$
  before the softmax, so their weights become exactly zero. Combined with the right shift,
  this guarantees that the prediction for position $i$ depends only on outputs before $i$.
\item \textbf{Encoder--decoder (cross) attention.} An extra sub-layer where the \emph{queries}
  come from the decoder and the \emph{keys and values} come from the encoder output. This is
  the classical seq2seq attention, now expressed in Q/K/V form: use the current partial
  translation as query to look up relevant source content.
\end{enumerate}

\paragraph{Final linear and softmax layer.} The decoder's top-layer vector is projected by a
linear layer to vocabulary size (the ``logits'' vector), then softmax converts it into a
probability distribution. Greedy decoding takes the argmax; beam search keeps several
hypotheses; sampling with temperature and top-$k$/top-$p$ is used for creative generation.

\paragraph{Counterfactual.} \emph{What if we removed the causal mask during training?} The
model would see the answer while predicting it, achieve near-zero training loss, and be
completely useless at generation time, when future tokens do not exist. The mask is not a
performance trick --- it is what makes the training objective match the inference procedure.

\subsection{Byte-Pair Encoding (BPE) and Subword Tokenization}

\paragraph{The problem.} Word-level tokenization needs a fixed vocabulary. Real language
defeats it:
\begin{itemize}
\item German compounds: \emph{Krankenhaus} = \emph{kranken} (sick) + \emph{haus} (house)
  = hospital. Every compound would be a new word.
\item English inflection: \emph{low}, \emph{lower}, \emph{lowest} become three unrelated
  entries.
\item Any unseen word maps to \texttt{<UNK>}, destroying information --- fatal for names,
  typos, code and morphologically rich languages.
\end{itemize}
Character-level tokenization avoids \texttt{<UNK>} but makes sequences very long and forces
the model to relearn spelling. Stemming and lemmatisation are not used, because they are
language-specific, lossy, and rule-based.

\paragraph{What BPE is.} A data-driven compression algorithm turned into a tokenizer. Start
from characters (or bytes) and repeatedly merge the most frequent adjacent pair into a new
symbol, for a fixed number of merges. Frequent words end up as single tokens; rare words
decompose into meaningful pieces.

\paragraph{Worked example (from the lecture).}
\[
\texttt{aaabdaaabac} \rightarrow \texttt{ZabdZabac}\ (Z{=}aa) \rightarrow
\texttt{ZYdZYac}\ (Y{=}ab) \rightarrow \texttt{XdXac}\ (X{=}ZY)
\]

\paragraph{Advantages.} Fixed, tunable vocabulary size; no out-of-vocabulary tokens (with
byte-level BPE, literally any input is representable); shares subwords across related words;
gives sensible splits for morphology.

\paragraph{Disadvantages.} Splits are statistical, not linguistic, so they sometimes cut
across morpheme boundaries; token counts are unfair across languages (non-English and
non-Latin scripts consume more tokens for the same content, making them more expensive and
effectively giving them less context); and it makes character-level tasks (counting letters,
rhyming) unnaturally hard for the model.

\paragraph{The family, and who uses what.}
\begin{center}
\begin{tabularx}{\textwidth}{lY}
\textbf{Scheme} & \textbf{Used by} \\
\hline
BPE / byte-level BPE & GPT-2, GPT-3, RoBERTa \\
WordPiece & BERT (merges by likelihood gain rather than raw frequency) \\
SentencePiece & T5, LLaMA (treats input as a raw byte stream, so no pre-tokenization is
needed and whitespace is reversible; works for languages without spaces) \\
\end{tabularx}
\end{center}

\subsection{Transformer: Summary Judgement}

\paragraph{Problems it solves.} Parallel training; constant-path-length long-range
dependencies; a single architecture that transfers to almost every modality.

\paragraph{What it cannot solve.}
\begin{itemize}
\item Quadratic cost in sequence length --- long documents remain expensive.
\item It has no memory beyond its context window; it is stateless between calls.
\item It is extremely data- and compute-hungry, because it encodes almost no inductive bias
  (no locality, no recurrence). On small datasets an LSTM or CNN often wins.
\item Fixed maximum length: once trained, extending context is non-trivial.
\end{itemize}

\paragraph{What it gave rise to.} The pretraining era (BERT, GPT), and with it a whole new
set of problems: models too large to train per task, too large to deploy, misaligned with
user intent, and confidently wrong. Parts IV--IX are the responses to those.

\paragraph{Use it when.} You have substantial data and GPUs, need long-range dependencies,
or want to build on a pretrained checkpoint (which today is nearly always).

\paragraph{Do not use it when.} Data is small and no suitable pretrained model exists;
sequences are extremely long and you cannot afford $O(n^2)$; or the deployment target is
tiny and streaming, where a GRU may be preferable.

\section{Part IV: Pretraining, BERT and GPT}

\subsection{Pretraining and Transfer Learning}

\paragraph{The motivating problem.} Training a Transformer from scratch is exceptionally
difficult: not enough hardware, not enough time, and rarely enough labelled data for the
specific task you care about.

\paragraph{The insight.} Most of what a model needs to know --- syntax, word senses, world
facts, discourse conventions --- is task-independent. So let an organisation with the
hardware train a general \emph{language representation} once on unlabelled text, and let
everyone else adapt that representation to their task with a small labelled dataset.

\paragraph{Two ways to use a pretrained model.}
\begin{itemize}
\item \textbf{Fine-tuning:} add a thin task-specific output layer and continue training all
  (or most) parameters on your labelled data. Usually more accurate.
\item \textbf{Feature extraction:} freeze the model, run your data through it once, and use
  the resulting contextual vectors as inputs to any classifier. Cheaper, and lets you
  precompute embeddings; useful when you cannot afford to fine-tune or must serve many tasks
  from one frozen backbone. For BERT, concatenating the last four layers was found to be the
  best feature-based configuration.
\end{itemize}

\paragraph{Why unlabelled text is enough.} Because the training signal is manufactured from
the text itself --- predict a hidden word, or predict the next word. This is
\emph{self-supervised} learning, and it is why the internet becomes a training set.

\subsection{BERT}
\label{sec:bert}

\paragraph{What it is.} Bidirectional Encoder Representations from Transformers: a stack of
Transformer \emph{encoders} pretrained to produce deep bidirectional representations of text
by jointly conditioning on both left and right context in all layers. A pretrained BERT can
be fine-tuned with just one additional output layer to reach state of the art on a wide
range of tasks.

\paragraph{The idea, plainly.} Earlier models read a sentence in one direction (or ran two
independent one-directional models and glued them together). BERT wanted a representation of
``bank'' that is informed by ``river'' on the left and ``flooded'' on the right
\emph{simultaneously, in every layer}. The obstacle described in the bidirectional RNN
section applies: a model that can see the whole sentence cannot be trained to predict the
next word, because it would see the answer. BERT's solution is to change the task.

\paragraph{Sizes.}
\begin{center}
\begin{tabularx}{\textwidth}{lYYY}
\textbf{Model} & \textbf{Encoder layers} & \textbf{Heads / hidden units} & \textbf{Parameters} \\
\hline
BERT-base & 12 & 12 heads, 768 hidden & 110 million \\
BERT-large & 24 & 16 heads, 1024 hidden & 340 million \\
\end{tabularx}
\end{center}
Tokenization is WordPiece. Inputs are built as
\texttt{[CLS] sentence A [SEP] sentence B [SEP]}, with token, segment and position
embeddings summed. The final hidden state of \texttt{[CLS]} is the aggregate sequence
representation used for classification.

\paragraph{Pretraining task 1: Masked Language Modelling (MLM).} Choose 15\% of tokens at
random as prediction targets. For each chosen token:
\begin{itemize}
\item 80\% of the time, replace it with \texttt{[MASK]};
\item 10\% of the time, replace it with a random token;
\item 10\% of the time, leave it unchanged.
\end{itemize}
The model must predict the original token. Because it does not know which positions were
corrupted or left alone, it is forced to build a good contextual representation of
\emph{every} token.

\paragraph{Why the 80/10/10 split?} \texttt{[MASK]} never appears during fine-tuning, so a
model trained on 100\% \texttt{[MASK]} would suffer a pretrain/fine-tune mismatch. The
random-token 10\% forces the model to keep a distributional representation of every input
token rather than trusting it blindly; the unchanged 10\% biases the representation toward
the actual observed word. The paper's ablation shows fine-tuning is fairly robust to the
exact ratios, but the feature-based setting (where representations are frozen and the
mismatch cannot be corrected) degrades noticeably with a 100\% \texttt{[MASK]} strategy.
That is the sharpest way to state the justification.

\paragraph{Pretraining task 2: Next Sentence Prediction (NSP).} Given sentences A and B,
predict whether B actually followed A. Half the training pairs are genuine successors and
half are random sentences from the corpus. Motivation: important tasks --- question
answering, natural language inference --- are about relationships \emph{between} sentences,
which single-sentence language modelling never teaches.

\paragraph{Advantages.}
\begin{itemize}
\item Deeply bidirectional context in every layer; excellent for understanding tasks.
\item One pretrained model, many downstream tasks, with only a small head added.
\item Dramatically reduces labelled-data requirements.
\end{itemize}

\paragraph{Disadvantages.}
\begin{itemize}
\item \textbf{Cannot generate text.} It is an encoder; there is no autoregressive decoder.
\item Only 15\% of tokens contribute to the loss per batch, so pretraining converges more
  slowly per step than left-to-right modelling. (The paper's ablation confirms MLM converges
  slightly slower but overtakes left-to-right in absolute accuracy almost immediately.)
\item Pretrain/fine-tune mismatch from the artificial \texttt{[MASK]} symbol.
\item Requires a fine-tuning step and labelled data per task --- no zero-shot ability.
\item Fixed 512-token limit; large, slow to serve.
\end{itemize}

\paragraph{What it cannot solve.} Open-ended generation, translation, summarisation, dialogue
--- anything that produces a sequence.

\paragraph{What it gave rise to.} RoBERTa (drop NSP, train longer on more data --- showing
NSP was not essential), ALBERT (parameter sharing), DistilBERT (distillation), ELECTRA
(replaced-token detection to use 100\% of tokens as signal), and the general recognition
that pretraining objectives are a design space in their own right.

\paragraph{Use it when.} Classification, sentiment, NER, extractive QA, sentence similarity,
retrieval embeddings --- any task where you need to \emph{understand} a fixed text and where
a 110M-parameter model that runs cheaply beats calling an LLM API.

\paragraph{Do not use it when.} You need generation, or zero/few-shot behaviour with no
labelled data, or inputs longer than 512 tokens.

\paragraph{Counterfactual 1.} \emph{What if BERT had been trained left-to-right instead of
with masking?} It would be GPT. You would gain generation ability and lose the right-context
information that makes it so strong at classification and extractive QA. The paper's own
ablation shows the left-to-right variant is worse on MNLI.

\paragraph{Counterfactual 2.} \emph{What if masking used 100\% \texttt{[MASK]}?} See above:
fine-tuning barely changes, but frozen feature quality drops (NER feature-based falls from
94.9 to 94.0 in the ablation), because the model never learns useful representations for
\emph{uncorrupted} tokens --- which are the only kind it will ever see downstream.

\paragraph{Counterfactual 3.} \emph{What if we masked 50\% of tokens instead of 15\%?}
Too little surviving context to reconstruct from, so representations would be poorer. Mask
too few, and each expensive forward pass yields almost no training signal. 15\% is a
compromise between signal per step and context per prediction.

\paragraph{Counterfactual 4.} \emph{What if NSP were dropped?} RoBERTa did exactly this and
performed \emph{better}, given more data and longer training. The lesson: NSP is a weak
objective, partly because ``random sentence from a different document'' can be detected on
topic alone rather than by genuine discourse reasoning.

\subsection{GPT: Decoder-Only Generative Models}

\paragraph{What it is.} A stack of Transformer \emph{decoder} blocks (masked self-attention
plus feed-forward; no cross-attention, since there is no encoder) trained with causal
language modelling: predict the next token.

\paragraph{The idea, plainly.} Nearly every NLP task can be written as text continuation.
``Translate to French: I am hungry $\rightarrow$'' is a prompt; the translation is simply the
most likely continuation. So instead of designing one architecture per task, train one model
to continue text extremely well and phrase every task as a continuation.

\paragraph{Why decoder-only.} Autoregressive next-token prediction requires that position $i$
cannot see position $i+1$; that is exactly what the decoder's masked self-attention provides.
The encoder--decoder split is only needed when input and output are genuinely different
sequences; if you concatenate prompt and answer into one stream, a single masked stack
suffices.

\paragraph{The generations.}
\begin{center}
\begin{tabularx}{\textwidth}{lY}
\textbf{Model} & \textbf{What changed} \\
\hline
GPT-1 (2018) & ``Improving Language Understanding by Generative Pre-Training.''
Established the recipe: generative pretraining on unlabelled text, then discriminative
fine-tuning per task. About 120 million parameters. \\
GPT-2 (2019) & A direct scale-up of GPT-1 --- roughly ten times the parameters (1.5 billion)
and ten times the data. Trained on WebText: pages linked from Reddit posts with at least 3
karma before December 2017, about 45 million links and 40 GB of text. The conceptual leap:
if a language model is powerful enough and the corpus diverse enough, the corpus itself
contains demonstrations of many tasks, so \emph{fine-tuning becomes unnecessary}. Hence the
title: ``Language Models are Unsupervised Multitask Learners.'' Last GPT model released
publicly. \\
GPT-3 (2020) & Scaled again to 175 billion parameters; roughly 300 billion training tokens
drawn from CommonCrawl, WebText, English Wikipedia and two book corpora (about 570 GB of
plaintext). Reported cost on the order of 355 GPU-years and several million dollars. Its
defining property is being a much better \emph{in-context learner}. API access only.
InstructGPT and GPT-3.5 added instruction tuning and RLHF, and made the current LLM boom
possible. \\
GPT-4 (2023) & Multimodality, heavy instruction tuning, reinforcement learning from human and
AI feedback, larger context windows (up to 32k tokens in the released variants) and much
stronger chat behaviour. \\
\end{tabularx}
\end{center}

\paragraph{The trajectory in one sentence.} GPT-1 = pretrain then fine-tune; GPT-2 = maybe you
do not need to fine-tune; GPT-3 = you definitely do not, just show examples in the prompt;
GPT-4 = now make it obedient, safe and multimodal.

\paragraph{Advantages.} Generation; one model for all tasks; zero- and few-shot ability;
capabilities emerge with scale that were not explicitly trained for.

\paragraph{Disadvantages.} Unidirectional context, so per-parameter it is weaker than BERT at
pure understanding tasks; extremely expensive to train and serve; hallucinates confidently;
outputs are stochastic and hard to guarantee.

\paragraph{What it cannot solve.} Guaranteed factual accuracy; knowledge after its training
cut-off; access to private data; genuine multi-step deliberation in a single forward pass;
arithmetic reliability.

\paragraph{What it gave rise to.} The whole second half of this syllabus: efficiency methods
(the models are too big), alignment methods (they are not obedient), RAG (they are not
grounded), agents (they are passive), and evaluation research (nobody knows how to score
them).

\paragraph{Use it when.} Generation, open-ended tasks, few-shot settings, or when you need one
system to cover many tasks.

\paragraph{Do not use it when.} You have a narrow, high-volume classification task with
labelled data --- fine-tuned BERT will be cheaper, faster and often more accurate --- or when
outputs must be verifiable and deterministic.

\paragraph{Counterfactual.} \emph{What if GPT had used an encoder--decoder instead?} You get
T5/BART. These are genuinely better for tasks with a clear input--output split (translation,
summarisation) because the encoder can be bidirectional over the input. But they are more
complex, and once prompting proved that everything can be expressed as continuation, the
simpler decoder-only stack won on scaling convenience.

\subsection{Choosing an Architecture: The Decision Table}

\begin{center}
\begin{tabularx}{\textwidth}{lYY}
\textbf{Family} & \textbf{Attention pattern} & \textbf{Best suited to} \\
\hline
Encoder-only (BERT, RoBERTa) & Bidirectional & Classification, NER, extractive QA,
sentence similarity, embeddings for retrieval \\
Decoder-only (GPT, LLaMA) & Causal (masked) & Free-form generation, chat, code, few-shot
tasks, anything open-ended \\
Encoder--decoder (original Transformer, T5, BART) & Bidirectional encoder + causal decoder
with cross-attention & Translation, summarisation, and other tasks with a distinct source
and target \\
\end{tabularx}
\end{center}

\paragraph{Rule of thumb for the exam.} \emph{Understand a fixed text} $\rightarrow$ encoder.
\emph{Produce free text} $\rightarrow$ decoder. \emph{Transform one text into another}
$\rightarrow$ encoder--decoder.

\section{Part V: How Large Language Models Are Pretrained}

\subsection{The Next-Token Prediction Objective}

\paragraph{The task.} Given \texttt{"The capital of Bangladesh is"}, produce
\texttt{"Dhaka"}. Formally, maximise the probability the model assigns to the correct next
token, over the whole corpus.

\paragraph{The loss.} Cross-entropy between the model's predicted distribution over the
vocabulary and the one-hot target. For a single position with true token $y$:
\[
\mathcal{L} = -\log p_\theta(y \mid x_{<t}).
\]
Perplexity, the standard reporting metric, is $\exp$ of the average of this loss.

\paragraph{Teacher forcing.} During training the model is fed the \emph{ground-truth}
previous tokens, never its own predictions. This is what allows every position in a sequence
to be trained in parallel in a single forward pass --- if the model had to consume its own
outputs during training, training would become sequential and lose the Transformer's main
advantage.

\paragraph{Autoregressive generation.} At inference there is no ground truth, so the model's
own output at step $t$ becomes its input at step $t+1$.

\paragraph{The tension between them (exposure bias).} The model is only ever trained on
perfect prefixes but is deployed on prefixes it wrote itself, which contain its own mistakes.
Once it makes an error it is in a state it never saw during training, and errors compound.
This is one root cause of degeneration and of hallucination cascades.

\paragraph{Counterfactual.} \emph{What if we trained without teacher forcing, always feeding
the model's own outputs?} Training would match inference (removing exposure bias), but early
in training the model's outputs are noise, so it would have almost nothing to learn from;
convergence would be far slower or fail entirely; and training would become sequential.
Scheduled sampling --- gradually mixing in the model's own predictions --- is the compromise,
but teacher forcing remains standard because parallelism matters more.

\subsection{Training Data}

\paragraph{Common sources.}
\begin{itemize}
\item \textbf{Common Crawl} --- petabyte-scale web pages. Huge but very noisy.
\item \textbf{Books} --- Project Gutenberg and other, less clearly licensed libraries. Long,
  coherent, high-quality prose; teaches long-range structure.
\item \textbf{Wikipedia} --- clean, structured, factual.
\item \textbf{Academic papers} --- ArXiv, PubMed. Technical vocabulary and reasoning.
\item \textbf{Code repositories} --- GitHub. Notably improves reasoning and structured output
  even for non-code tasks.
\end{itemize}

\paragraph{The governing principle.} Data quality outweighs raw quantity. Aggressive
filtering, deduplication and quality classifiers routinely beat simply adding more raw web
text.

\paragraph{Problems with the data.}
\begin{itemize}
\item \textbf{Quality:} spam, machine-generated text, boilerplate.
\item \textbf{Diversity:} heavy skew toward English and toward the demographics that write
  most on the web; low-resource languages are badly underserved.
\item \textbf{Contamination:} benchmark test sets appear in the crawl, so evaluation scores
  become unreliable (see Section~\ref{sec:eval}).
\item \textbf{Bias and toxicity:} the model inherits whatever the corpus contains.
\item \textbf{Legal and ethical status:} copyright, licensing, consent, and personal data
  scraped without permission.
\end{itemize}

\subsection{Scaling Laws and the Chinchilla Result}

\paragraph{What scaling laws are.} Empirical power-law relationships showing that loss falls
predictably as model parameters, dataset size and compute increase. They let you forecast the
performance of a model you have not yet trained, which is what makes multi-million-dollar
training runs financially defensible.

\paragraph{The Chinchilla correction (Hoffmann et al., 2022).} Before this paper, the field
scaled parameters far faster than data. The finding:
\begin{itemize}
\item Most existing models were \textbf{compute-inefficient}: too large for the amount of
  data they were trained on.
\item Model size and dataset size should be scaled \textbf{equally}. Given a fixed compute
  budget, performance improves more by balancing the two than by inflating the model.
\item \textbf{Doubling the training data gives comparable or better gains than doubling model
  size.}
\item \textbf{Overparameterised} models --- too many parameters relative to their training
  data --- underperform smaller models trained on more data.
\end{itemize}

\paragraph{The evidence.} Chinchilla, at 70 billion parameters, outperformed Gopher, a
280-billion-parameter predecessor, using the same compute budget spent on more tokens. The
practical heuristic that followed is roughly 20 training tokens per parameter.

\paragraph{Why this matters beyond training.} A smaller model trained on more data is cheaper
to \emph{serve}, forever. Training happens once; inference happens billions of times. This is
the economic reason the industry moved toward smaller, longer-trained models --- and directly
why LLaMA (7B/13B trained on trillions of tokens) was possible.

\paragraph{What scaling laws cannot tell you.} Which capabilities will appear and when;
whether the model will be safe, honest or useful; or what happens when high-quality human
text runs out.

\paragraph{Counterfactual.} \emph{What if the field had kept scaling parameters only, as
before Chinchilla?} We would have a small number of enormous, undertrained, extremely
expensive-to-serve models, and no 7B-parameter models good enough to run on a laptop. The
entire open-weights ecosystem depends on this result.

\section{Part VI: LLaMA and Modern Architectural Refinements}

\subsection{The LLaMA Family}

\paragraph{What it is.} A family of decoder-only LLMs published by Meta AI --- arguably the
most popular freely available LLM --- spanning several generations and parameter sizes
(7B, 13B, 34B, 70B and others).

\paragraph{Why it matters.} It demonstrated the Chinchilla lesson in public: comparatively
small models trained on very large token counts, released with open weights, so that
academics and small companies could fine-tune state-of-the-art models themselves. It created
the open LLM ecosystem.

\paragraph{Scale, to give a sense of cost.} LLaMA 2 70B was among the largest LLMs at
release: trained on roughly 10 TB of text, requiring about 6000 GPUs for 12 days, on the
order of $10^{24}$ floating-point operations, at a cost around 2 million US dollars.

\paragraph{Architecture at a glance.} LLaMA is the original Transformer decoder with four
substitutions. Compared side by side with the 2017 architecture:
\begin{center}
\begin{tabularx}{\textwidth}{lYY}
\textbf{Component} & \textbf{Original Transformer} & \textbf{LLaMA} \\
\hline
Normalisation & LayerNorm, applied after the residual add (post-norm) &
RMSNorm, applied to the sub-layer input (pre-norm) \\
Position information & Sinusoidal encoding added once to the input embeddings &
Rotary positional encoding (RoPE) applied to Q and K inside every attention layer \\
Attention & Multi-head attention & Grouped-query attention with a KV cache
(in the larger LLaMA 2 variants) \\
Feed-forward & ReLU MLP & SwiGLU gated MLP \\
\end{tabularx}
\end{center}

\subsection{Normalisation, and RMSNorm}

\paragraph{Why normalise at all?} As signals pass through many layers, their scale drifts;
activations can grow or shrink until gradients explode or vanish. Normalisation rescales
each layer's activations into a stable range, which permits higher learning rates, reduces
sensitivity to initialisation, smooths the loss landscape, and shortens training time
considerably.

\paragraph{LayerNorm.} Subtract the mean and divide by the standard deviation across the
feature dimension, then apply learned gain $\gamma$ and bias $\beta$:
\[
\mathrm{LN}(x) = \gamma \odot \frac{x - \mu}{\sqrt{\sigma^2 + \epsilon}} + \beta .
\]

\paragraph{RMSNorm.} Drop the mean subtraction and the bias; rescale by the root mean square
only:
\[
\mathrm{RMSNorm}(x) = \gamma \odot \frac{x}{\sqrt{\frac{1}{d}\sum_{i=1}^{d} x_i^2 + \epsilon}} .
\]

\paragraph{Why the simplification works.} Empirically, the benefit of LayerNorm comes almost
entirely from \emph{rescaling}, not from \emph{recentring}. Removing the mean computation
removes one pass over the vector and several operations per token, giving a measurable
speed-up (commonly cited around 7--15\% of training time) at essentially no quality cost. At
LLaMA's scale, a few percent is worth many GPU-days.

\paragraph{Counterfactual.} \emph{What if LLaMA had kept post-norm LayerNorm?} Training very
deep stacks would require careful learning-rate warm-up and would be more prone to
divergence, and each step would be slightly slower. Pre-norm keeps an unobstructed residual
path from input to output, which is what makes 80-layer models trainable.

\subsection{Positional Encoding, Revisited: Absolute, Relative and Rotary}

\paragraph{Absolute (original Transformer).} Add a vector that encodes ``I am token 7''. The
model must infer relative distance from the difference of two absolute encodings. It also
does not extrapolate well past the trained length.

\paragraph{Relative (Shaw et al.).} Inject a learned embedding of the \emph{offset} $i - j$
directly into the attention score computation. Linguistically this is the right quantity ---
what matters is that the adjective is two words before the noun, not that it is at absolute
position 41. The drawback is extra computation and memory inside every attention operation.

\paragraph{Rotary positional encoding (RoPE), used by LLaMA.} Do not \emph{add} anything.
\emph{Rotate} the query and key vectors by an angle proportional to their position, treating
consecutive pairs of dimensions as points in a 2-D plane:
\[
q_m' = R_m q_m, \qquad k_n' = R_n k_n, \qquad
q_m'^\top k_n' = q_m^\top R_{n-m} k_n .
\]
The crucial property: because rotations compose, the dot product of a rotated query and a
rotated key depends only on the \emph{difference} $n - m$. So an absolute-style operation
produces relative-position behaviour for free.

\paragraph{Why it is the modern default.}
\begin{itemize}
\item Gives relative position information at close to zero extra cost, with no added
  parameters.
\item Attention naturally decays with distance, a useful inductive bias.
\item Applied to Q and K at every layer, not just once at the input, so position information
  does not fade with depth.
\item Extends to longer contexts by interpolating or rescaling the rotation frequencies ---
  the standard trick behind long-context fine-tuning.
\end{itemize}

\paragraph{Counterfactual.} \emph{What if RoPE rotated the value vectors as well?} The output
would become position-dependent in an uncontrolled way and the clean relative-position
identity above would break. Rotation is deliberately confined to the matching operation
(Q and K), not the content being retrieved (V).

\subsection{Multi-Query and Grouped-Query Attention}

\paragraph{The problem it addresses.} Modern GPUs compute far faster than they can move data.
Arithmetic is cheap; reading memory is expensive. So the bottleneck in \emph{inference} is
often memory bandwidth, not FLOPs.

\paragraph{Where the memory goes: the KV cache.} During autoregressive generation, the keys
and values of all previous tokens are cached so they need not be recomputed for each new
token. With multi-head attention, each of $h$ heads stores its own K and V for every position.
For a long conversation this cache becomes gigabytes, and it must be read from memory for
every single generated token.

\paragraph{The spectrum of solutions.}
\begin{center}
\begin{tabularx}{\textwidth}{lYY}
\textbf{Scheme} & \textbf{Key/value heads} & \textbf{Effect} \\
\hline
Multi-head attention (MHA) & One K/V pair per query head & Best quality, largest cache \\
Multi-query attention (MQA) & A single K/V pair shared by all query heads & Smallest cache,
fastest, some quality loss \\
Grouped-query attention (GQA) & Query heads are split into groups; one K/V pair per group &
Near-MHA quality at close to MQA speed \\
\end{tabularx}
\end{center}
Queries stay fully multi-headed in all three; only the keys and values are shared.

\paragraph{Why it works.} The expressive power of attention lives mostly in having many
different \emph{queries} --- many different questions being asked. Having many distinct
\emph{keys and values} matters less. GQA exploits that asymmetry.

\paragraph{Important correction.} The lecture slides list grouped multi-query attention among
``LLaMA 1 new concepts''. In the published papers, LLaMA 1 used standard multi-head
attention; GQA was introduced in \textbf{LLaMA 2}, and only in the 34B and 70B variants
(7B and 13B kept MHA). Know both the slide claim and the correct attribution.

\paragraph{Counterfactual.} \emph{What if we shared the queries instead of the keys and
values?} All heads would ask the same question of different content, which destroys the
entire point of multi-head attention, and it would not shrink the KV cache at all --- queries
are not cached, since they are recomputed for the current token only.

\subsection{SwiGLU}

\paragraph{What it is.} A gated feed-forward layer that replaces the ReLU MLP:
\[
\mathrm{SwiGLU}(x) = \big(\mathrm{Swish}(xW_1) \odot xW_3\big)W_2,
\qquad \mathrm{Swish}(z) = z \cdot \sigma(z).
\]

\paragraph{The idea, plainly.} A ReLU MLP applies one nonlinearity and passes the result on.
SwiGLU computes two projections: one becomes a smooth \emph{gate} and the other the
\emph{content}, and they are multiplied elementwise. It is the same gating intuition as the
LSTM --- let the network decide, per feature and per token, how much signal to let through
--- applied inside the feed-forward block.

\paragraph{Practical details.} It uses three weight matrices instead of two, so to keep the
parameter count unchanged the hidden dimension is reduced to about $\frac{2}{3}$ of the usual
$4 d_{\text{model}}$. Swish is smooth and non-monotonic, unlike ReLU, which avoids the ``dead
neuron'' problem where a unit stuck in the negative region receives zero gradient forever.

\paragraph{Honest caveat.} The original paper introducing GLU variants attributes their
success to ``divine benevolence'' --- that is, they work reliably in practice but there is no
crisp theoretical reason. Saying this in an exam is better than inventing a justification.

\subsection{Long Context Models}

\paragraph{Definition.} The context window (or context length) of an LLM is the amount of
text, in tokens, that the model can consider or ``remember'' at any one time.

\paragraph{The trajectory.} 4K (GPT-3) $\rightarrow$ 128K (GPT-4 Turbo) $\rightarrow$ 1M+
(Gemini 1.5).

\paragraph{Why it matters.} Whole-book analysis, entire codebases, long legal contracts,
extended multi-turn conversations, and RAG pipelines that can afford to retrieve generously.

\paragraph{Limitations.}
\begin{itemize}
\item \textbf{``Lost in the middle.''} Retrieval accuracy is high for material at the
  beginning and end of the context and drops sharply for material in the middle. A large
  window does not guarantee the model uses all of it.
\item \textbf{Quadratic attention cost.} Doubling context quadruples attention compute and
  memory. FlashAttention reduces the memory cost by never materialising the full $n \times n$
  matrix (not covered in this course), but the asymptotic compute cost remains.
\item Cost and latency scale with context, so a long prompt is expensive on every call.
\item Quality within the window can degrade even when the model does not error.
\end{itemize}

\paragraph{Use it when.} The task genuinely requires global reasoning over a document.
\paragraph{Do not use it when.} You only need a few relevant passages --- RAG is cheaper,
more current and more auditable than stuffing everything into context.

\subsection{Mixture of Experts (MoE)}

\paragraph{Core idea.} Replace the single feed-forward block in each layer with many parallel
``expert'' feed-forward sub-networks, and use a small learned \emph{router} to send each token
to only a few of them. Most parameters exist but stay dormant for any given token.

\paragraph{The canonical example.} Mixtral 8x7B: 8 experts, each about 7B, totalling 47B
parameters --- not 56B, because only the feed-forward layers are expert-specific; token
embeddings, attention layers and so on are shared across experts. Only 2 experts are active
per token, so roughly 13B parameters are used per forward pass. Other examples: DeepSeek-V3
(256 experts, 8 active) and Qwen 2.5-MoE.

\paragraph{The trade it makes.} Decouples \emph{capacity} (total parameters, i.e. how much the
model can know) from \emph{compute} (active parameters, i.e. how much it costs to run a
token). You get the knowledge of a 47B model at roughly the inference cost of a 13B model.

\paragraph{Advantages.} Far better quality per FLOP; faster training and inference for a given
knowledge capacity; natural specialisation among experts.

\paragraph{Disadvantages.}
\begin{itemize}
\item \textbf{All} parameters must still fit in memory, even the inactive ones --- so VRAM
  requirements follow the 47B, not the 13B.
\item Training is unstable: routers collapse toward a few favourite experts unless an
  auxiliary load-balancing loss is added.
\item Complex to serve efficiently, especially across multiple GPUs.
\item Fine-tuning is harder and more prone to overfitting.
\end{itemize}

\paragraph{Use it when.} You are compute-bound at inference but have plenty of memory, and you
want maximum capability per FLOP.
\paragraph{Do not use it when.} Memory is the binding constraint (edge devices), or you need a
simple, easily fine-tuned model.

\paragraph{Counterfactual.} \emph{What if the router were random instead of learned?} Experts
would never specialise; you would effectively be training an ensemble of identical average
models and the capacity advantage would vanish. Conversely, \emph{what if we removed the
load-balancing loss?} The router would send nearly all tokens to a handful of experts, the
rest would be starved of gradient, and the effective model would shrink to those few ---
``expert collapse''.

\section{Part VII: Faster, Cheaper, Deployable Models}

The four techniques in this part answer one question: the pretrained model is excellent but
too large to fine-tune or serve. What can be removed?

\begin{center}
\begin{tabularx}{\textwidth}{lYY}
\textbf{Technique} & \textbf{What it reduces} & \textbf{What stays the same} \\
\hline
Quantization & Bits per weight & Architecture, parameter count \\
PEFT / LoRA & Trainable parameters & Model size at inference \\
Pruning & Number of weights or structures & Precision of what remains \\
Distillation & The whole model (new, smaller one) & Nothing --- it is a different model \\
\end{tabularx}
\end{center}

\subsection{Quantization}

\paragraph{What it is.} A technique to reduce the computational and memory cost of inference
by representing weights and activations with low-precision data types, such as 8-bit integers
(int8), instead of the usual 32-bit floating point (float32).

\paragraph{Why fewer bits helps.} The model needs less memory storage; it consumes less energy
(in theory); and operations such as matrix multiplication run much faster with integer
arithmetic. In practice, memory bandwidth is the main win: a 4-bit model moves one eighth as
many bytes per token as a float32 model.

\paragraph{Two regimes.}
\begin{itemize}
\item \textbf{Post-training quantization (PTQ).} Quantize an already-trained model. No
  retraining, so it is fast and cheap; accuracy may drop.
\item \textbf{Quantization-aware training (QAT).} Simulate quantization during training ---
  the forward pass uses quantized values while the backward pass keeps full precision.
  Recovers more accuracy, but requires a training run.
\end{itemize}

\paragraph{Common data types and accumulation types.}
\begin{center}
\begin{tabularx}{\textwidth}{lY}
\textbf{Storage type} & \textbf{Accumulation type} \\
\hline
float16 & float16 \\
bfloat16 & float32 \\
int16 & int32 \\
int8 & int32 \\
\end{tabularx}
\end{center}
The \textbf{accumulation data type} specifies the type of the result of accumulating (adding,
multiplying) values of the storage type. Example: two int8 values $A = 127$ and $B = 127$;
their sum $C = 254$ far exceeds the largest representable int8 value. Without a wider
accumulator, the result would overflow and the precision loss would make the whole
quantization pointless.

\paragraph{Case 1: float32 $\rightarrow$ float16.} Straightforward, because both types follow
the same representation scheme. Two checks are required:
\begin{itemize}
\item \emph{Does the hardware support it?} Intel CPUs long supported float16 only as a
  storage type, converting to float32 for computation; full support arrived with Cooper Lake
  and Sapphire Rapids.
\item \emph{Is the operation sensitive to lower precision?} The $\epsilon$ in LayerNorm is
  typically about $10^{-12}$, but the smallest representable float16 value is about
  $6 \times 10^{-5}$; the small value underflows to zero and can produce NaNs. The same risk
  exists at the top of the range for large values. This is exactly why bfloat16 --- which
  keeps float32's exponent range and sacrifices mantissa bits instead --- became the standard
  training format.
\end{itemize}

\paragraph{Case 2: float32 $\rightarrow$ int8.} Only 256 distinct values are representable in
int8, while float32 spans an enormous range. The task is to project a float range $[a,b]$ onto
the int8 grid. The \textbf{affine quantization scheme} is
\[
x = S\,(x_q - Z), \qquad x_q = \mathrm{round}\!\left(\frac{x}{S} + Z\right),
\]
where $x_q$ is the int8 value corresponding to $x$; $S$ is the \emph{scale}, a positive
float32; and $Z$ is the \emph{zero-point}, the int8 value that maps exactly to float 0.
Representing zero exactly matters because zero appears everywhere in machine learning
(padding masks, ReLU outputs, sparse weights); if it were only approximated, error would be
injected into every padded position.

\paragraph{Advantages.} Large reductions in memory and bandwidth; often 2--4$\times$ faster
inference; makes large models runnable on consumer hardware; requires no architectural change.

\paragraph{Disadvantages.} Some accuracy loss, which grows sharply below about 4 bits;
outlier activation values in LLMs are notoriously hard to quantize; hardware support is
uneven; and quantized models can be harder to fine-tune further.

\paragraph{Use it when.} Deploying under memory or latency constraints, or serving many
concurrent users.
\paragraph{Do not use it when.} You are still training or exploring; or accuracy is critical
and unverified; or your hardware lacks native low-precision kernels, in which case you may get
the memory saving without the speed-up.

\paragraph{Counterfactual.} \emph{What if we quantized with a scale but no zero-point
(symmetric quantization)?} The mapping is simpler and slightly faster, and it works well for
weights, which are roughly symmetric about zero. It works badly for post-ReLU activations,
which are all non-negative --- half the int8 range would be wasted on values that never occur.

\subsection{Parameter-Efficient Fine-Tuning (PEFT)}

\paragraph{What it is.} A family of methods that freeze most of the pretrained model's
parameters and layers, and train only a small number of added task-specific parameters. By
training a small set of parameters and preserving most of the pretrained structure, PEFT saves
time and computational resources.

\paragraph{Common techniques.}
\begin{itemize}
\item \textbf{Adapters.} Insert small bottleneck feed-forward layers between the existing
  Transformer layers and train only those.
\item \textbf{LoRA (Low-Rank Adaptation).} Do not modify $W$; learn a low-rank correction
  $\Delta W = BA$, where $B \in \mathbb{R}^{d \times r}$ and $A \in \mathbb{R}^{r \times k}$
  with $r \ll d$. The forward pass becomes $h = Wx + BAx$. Training a $4096 \times 4096$ layer
  with $r = 8$ means training about 65{,}000 parameters instead of 16.8 million.
\item Prompt/prefix tuning: learn continuous ``virtual tokens'' prepended to the input.
\end{itemize}

\paragraph{Why low rank is enough.} The hypothesis is that adapting a model to a new task
requires only a low-dimensional change --- the pretrained model already contains the
capability, and fine-tuning mostly re-weights what it already knows.

\paragraph{Advantages (as listed in the course).} Increased efficiency; faster time-to-value;
no catastrophic forgetting; lower risk of overfitting; lower data demands; more accessible AI;
more flexible AI.

Expanding two of these: \emph{no catastrophic forgetting}, because the base weights are frozen
and cannot be overwritten; \emph{more flexible}, because LoRA adapters are tiny files
(megabytes) that can be swapped at serving time, so one base model can serve dozens of tasks.
With LoRA the adapter can also be merged back into $W$ after training, adding zero inference
latency.

\paragraph{Disadvantages.} Usually slightly below full fine-tuning on tasks requiring large
behavioural change; the rank $r$ is another hyperparameter; and unmerged adapters add a small
serving overhead. Crucially, PEFT does not make the model \emph{smaller} at inference --- the
full base model must still be loaded.

\paragraph{What it cannot solve.} Teaching the model genuinely new knowledge or a new language
it never saw in pretraining; PEFT adapts, it does not extend. It also does not help when the
base model is simply too large for your hardware.

\paragraph{Use it when.} Adapting a large pretrained model to a task or style with limited
data and one GPU; maintaining many task variants of a single base model.
\paragraph{Do not use it when.} You need to teach substantial new knowledge (consider
continued pretraining or RAG); or your model is small enough to fine-tune fully.

\paragraph{Counterfactual.} \emph{What if we set the LoRA rank $r$ equal to the full
dimension?} $\Delta W$ becomes an unconstrained full-rank matrix and you have simply
reimplemented full fine-tuning, with all its memory cost and its risk of catastrophic
forgetting. The constraint \emph{is} the method.

\subsection{Pruning}

\paragraph{The procedure.}
\begin{enumerate}
\item Train the model.
\item Prune --- for example remove 30\% of the network.
\item Fine-tune to recover accuracy.
\end{enumerate}

\paragraph{Two kinds.}
\begin{itemize}
\item \textbf{Magnitude (unstructured) pruning:} remove individual weights with the smallest
  absolute values, on the assumption that small weights contribute least. Produces sparse
  weight matrices with high theoretical compression, but generic hardware cannot exploit
  irregular sparsity, so the real speed-up is often disappointing.
\item \textbf{Structured pruning:} remove whole neurons, attention heads or layers. Lower
  compression per unit of accuracy lost, but the result is a genuinely smaller dense model,
  which is better for hardware and gives real speed-ups.
\end{itemize}

\paragraph{The practical difficulty flagged in the lecture.} Recovery fine-tuning is
necessary, but PEFT may not be strong enough to repair the damage of removing 30\% of the
network, while a full fine-tune of an LLM is expensive. Pruning large models is therefore
harder than it sounds, and this is one reason distillation is often preferred at LLM scale.

\paragraph{Counterfactual.} \emph{What if we pruned without the recovery fine-tune?} Accuracy
would fall sharply. Removing weights disturbs the balance of every downstream layer; the
remaining weights must be readjusted to compensate. Pruning is best understood as
``damage then repair'', not simply ``delete''.

\subsection{Knowledge Distillation}

\paragraph{What it is.} Train a small ``student'' model to mimic a large ``teacher'' model.

\paragraph{How it works.}
\begin{enumerate}
\item The teacher LLM generates \emph{soft labels} --- full probability distributions --- over
  a dataset.
\item The student is trained to match the teacher's output distribution, minimising the KL
  divergence between the two distributions.
\end{enumerate}

\paragraph{Why soft labels beat hard labels.} A hard label says only ``this is a cat''. A soft
distribution says ``85\% cat, 12\% dog, 0.01\% aeroplane'', which encodes the teacher's
learned similarity structure --- often called \emph{dark knowledge}. The student learns not
just the answer but the shape of the teacher's uncertainty, so it learns far more per example
than it would from the hard labels alone.

\paragraph{Examples.} DistilBERT: 40\% smaller, retaining 97\% of BERT's accuracy.
TinyLlama (1.1B) trained from LLaMA 2 outputs.

\paragraph{Best for.} Creating tiny, fast, specialised models --- for instance an on-device
chat assistant.

\paragraph{Disadvantages.} Requires running the large teacher over a large dataset (expensive);
the student inherits the teacher's biases and errors; the student generally cannot exceed the
teacher; and licensing may forbid distilling from a commercial API.

\paragraph{Counterfactual.} \emph{What if the student were trained on the teacher's hard
argmax outputs instead of soft distributions?} You would lose the dark knowledge and be back
to ordinary supervised training on machine-generated labels. Distillation would still help a
little (the teacher's labels are cleaner than raw web data) but the characteristic efficiency
of distillation comes from the soft targets.

\section{Part VIII: Instruction Tuning and Alignment}

\subsection{Pretrained Models Are Not Chatbots}

A pretrained LLM is a text \emph{continuer}, not an \emph{assistant}. Ask it ``What is the
capital of Bangladesh?'' and a plausible continuation is another exam question, because that
is what usually follows a question on the web. It is not disobedient; it was never trained to
obey. The remaining stages fix that.

\begin{center}
\begin{tabularx}{\textwidth}{lYY}
\textbf{Stage} & \textbf{What it teaches} & \textbf{Data required} \\
\hline
Pretraining & Language, facts, reasoning patterns & Trillions of unlabelled tokens \\
Instruction tuning / SFT & Format: how to follow an instruction & Thousands to millions of
(instruction, response) pairs \\
RLHF / DPO & Preference: which of two valid answers humans prefer & Human comparisons of
model outputs \\
Constitutional AI & Principles, applied by the model to itself & A written constitution, plus
model-generated critiques \\
\end{tabularx}
\end{center}

\subsection{Instruction Tuning}

\paragraph{What it is.} Fine-tuning an LLM on a labelled dataset of instructional prompts and
their corresponding outputs, so that it responds to task-specific prompts phrased in natural
language. It is ordinary supervised learning on (input, output) pairs.

\paragraph{Example.}
\begin{itemize}
\item Instruction: What is the sentiment of this review?
\item Input: The movie was boring and too long.
\item Output: Negative
\end{itemize}

\paragraph{Why it works so well.} Training on \emph{many} tasks phrased as instructions
teaches the general skill ``an instruction is a thing to be obeyed'', which then transfers to
instructions never seen in training. FLAN (2021) demonstrated this cross-task generalisation.

\paragraph{Result.} The model becomes more user-friendly, more general-purpose, and better
aligned with human intent --- without changing the architecture at all.

\subsection{Alignment through Supervised Fine-Tuning}

\paragraph{Alignment, defined.} The process of shaping a model's behaviour to be helpful,
harmless and honest according to human intentions and values.

\paragraph{SFT is the easiest alignment method.}
\begin{enumerate}
\item Collect human demonstrations: (prompt, ideal response) pairs.
\item Fine-tune with standard cross-entropy loss on the responses.
\item Evaluate on held-out prompts.
\end{enumerate}

\paragraph{Its ceiling.} SFT teaches \emph{format}, not \emph{preferences}. It can show the
model what a good answer looks like, but it cannot express that answer A is 30\% better than
answer B, nor navigate the helpful-versus-harmless trade-off. It is also limited by the
quality of human writers --- the model imitates demonstrations rather than exceeding them.
And every demonstration must be written by a person, which is expensive.

\subsection{Reinforcement Learning from Human Feedback (RLHF)}

\paragraph{The pipeline.}
\begin{enumerate}
\item \textbf{Collect human preferences.} Show annotators two model responses, A and B, and
  ask which is better. Comparison is far easier and more reliable for humans than writing an
  ideal answer or assigning an absolute score.
\item \textbf{Train a reward model} that predicts the human preference score for a response.
  This turns sparse human judgement into a dense, automatic signal that can be queried
  millions of times.
\item \textbf{Optimise the policy} (the LLM) against the reward model using PPO (Proximal
  Policy Optimization).
\end{enumerate}

\paragraph{Why reinforcement learning at all?} Because ``be helpful'' has no ground-truth
target string to imitate. There is no single correct answer to copy, only better and worse
ones --- which is exactly the setting RL is designed for.

\paragraph{Assessment.} RLHF is expensive but produces the most helpful models (ChatGPT,
Claude). However, PPO is notoriously unstable; Direct Preference Optimization (DPO) is more
practical, because it optimises the same preference objective directly with a supervised-style
loss, dispensing with the separate reward model and the RL loop.

\paragraph{Disadvantages and failure modes.} Very expensive human annotation; annotator
disagreement and bias baked in; \emph{reward hacking}, where the model finds outputs the
reward model scores highly but humans dislike (excessive length, flattery, hedging); and an
``alignment tax'' --- some raw capability is often lost.

\paragraph{Counterfactual.} \emph{What if we optimised the reward model without a KL penalty
keeping the policy near the SFT model?} The policy would drift far from natural language and
exploit the reward model's blind spots, producing degenerate text that scores highly and reads
as nonsense. The KL constraint is what keeps RLHF anchored to fluent output.

\subsection{Constitutional AI}

\paragraph{What it is.} A method developed by Anthropic to train and align models to be
helpful, honest and harmless --- the three pillars of alignment --- without relying on humans
to manually review and score thousands of toxic or dangerous outputs.

\paragraph{How it works.} The model is given a written ``constitution'': a set of principles,
rules and guidelines. It then uses that constitution to critique and revise its own outputs.
The loop runs: constitution $\rightarrow$ red-teaming $\rightarrow$ self-critique $\rightarrow$
self-revision $\rightarrow$ supervised or reinforcement learning from the resulting feedback.

\paragraph{Advantages.} Dramatically reduces dependency on human-generated data; scales far
better; makes the values \emph{explicit and auditable} as a written document rather than
implicit in thousands of annotator judgements; and spares human workers from repeated exposure
to harmful content.

\paragraph{Disadvantages.} Someone must still write the constitution, so the value judgements
merely move upstream; the model may misinterpret abstract principles; and self-critique is
bounded by the model's own understanding --- it cannot catch a flaw it cannot recognise.

\section{Part IX: Using LLMs}

\subsection{In-Context Learning}

\paragraph{What it is.} Learning from examples provided in the prompt, with no weight updates.
Provide one example (one-shot) or a few (few-shot). Few-shot often beats zero-shot by a
significant margin.

\paragraph{Advantages.} Zero training cost; instant adaptability; no ML infrastructure needed;
you can change the ``task'' between two consecutive API calls.

\paragraph{Why does it work at all?} Nothing is being learned in the usual sense --- no
parameter changes. The leading explanation is that a sufficiently diverse pretraining corpus
contains countless implicit task demonstrations, so the model has learned to recognise and
continue patterns; the prompt's examples \emph{locate} an ability the model already has rather
than installing a new one. A related view is that the forward pass through attention performs
something functionally similar to a small gradient-descent-like update inside the activations.
Treat both as hypotheses, not settled fact.

\paragraph{Disadvantages.} Examples consume context on every call, so it is expensive at scale;
performance is sensitive to example choice, order and formatting; and it cannot match
fine-tuning when thousands of labelled examples exist.

\paragraph{Use it when.} Prototyping, few labelled examples, rapidly changing requirements.
\paragraph{Do not use it when.} High volume with stable requirements and abundant labels ---
fine-tune or distil instead, and stop paying for the examples on every request.

\subsection{Prompt Design}

\begin{itemize}
\item \textbf{Role prompting:} ``You are a professional data scientist\ldots''
\item \textbf{Structured prompting:} ``Step 1: Analyze. Step 2: Propose. Step 3: Conclude.''
\item \textbf{Delimiters:} use \verb|###| or triple quotes to separate instructions from input.
  This also mitigates prompt injection, since the model can tell data from commands.
\item \textbf{Output formatting:} ``Respond with JSON: \{"answer": ...\}''
\item \textbf{In-context examples:} one-shot or few-shot.
\item \textbf{Negative constraints:} ``Do not include any conversational filler\ldots''
\item \textbf{Escape hatch:} explicitly permit ``I don't know''. Without it the model's
  training pressure to produce a fluent continuation encourages fabrication; giving it a
  licensed way to decline reduces hallucination.
\end{itemize}

\subsection{Chain-of-Thought (CoT) Prompting}

\paragraph{What it is.} Asking the model not only to answer but also to generate a rationale
--- the reasoning steps by which it arrived at the answer.

\paragraph{Effect.} CoT significantly enhances performance on arithmetic, symbolic and other
logical reasoning tasks, including in the zero-shot case (``Let's think step by step'').

\paragraph{Why does it work?} Two complementary reasons, and the exam will want both:
\begin{enumerate}
\item \textbf{Computational.} A Transformer performs a fixed amount of computation per
  generated token. A hard problem may need more computation than one forward pass provides.
  Writing intermediate steps lets the model spend many forward passes on one problem --- the
  generated tokens act as an external working memory.
\item \textbf{Distributional.} Correct reasoning chains in the training corpus are usually
  followed by correct answers. Conditioning on a well-formed chain of reasoning shifts the
  model into a region of its distribution where correct answers are far more likely.
\end{enumerate}

\paragraph{Limitations.} It costs many extra tokens (latency and money); the stated rationale
is not guaranteed to be the actual cause of the answer, so it is a poor faithfulness guarantee
--- a model can produce a flawless-looking chain and a wrong answer, or the reverse; it mainly
helps large models, and can hurt small ones; and it does not help on tasks with no decomposable
structure.

\subsection{The Problem with Pure LLMs}

\begin{itemize}
\item \textbf{Hallucinations} --- fluent, confident, false statements.
\item \textbf{Outdated knowledge} --- frozen at the training cut-off.
\item \textbf{No access to private data} --- it never saw your company's documents.
\item \textbf{Expensive to retrain} --- you cannot re-pretrain a model every time a fact
  changes.
\end{itemize}
Solution: do not retrain. Give the LLM a retriever.

\subsection{Retrieval-Augmented Generation (RAG)}

\paragraph{What it is.} An architecture combining a \textbf{retriever}, which finds relevant
documents in a knowledge base, with a \textbf{generator}, which produces an answer conditioned
on those documents.

\paragraph{The retriever (the ``search engine'').}
\begin{itemize}
\item Convert documents into vector embeddings.
\item Store them in a vector database (Pinecone, Weaviate, FAISS).
\item At query time, embed the question and find the top-$k$ most similar documents.
\end{itemize}
Vector search matters because it matches on \emph{meaning}: a query about ``annual leave''
retrieves a passage about ``vacation policy'' even with no shared keywords.

\paragraph{The generator (the LLM).}
\begin{itemize}
\item Takes [retrieved documents] + [user query] + [instructions].
\item Generates an answer conditioned only on the retrieved context.
\item Can be any LLM. Notably, small LLMs (around 7B) work surprisingly well as generators
  when retrieval is good --- because the hard part becomes reading, not recalling.
\end{itemize}

\paragraph{Use cases.} Open-domain question answering; chatbots that must be factually
accurate; enterprise Q\&A over internal documents; keeping responses current without
retraining.

\paragraph{Advantages.} Knowledge updates instantly by editing the document store; sources can
be cited, so answers are auditable; private data never enters model weights; and it works with
a smaller, cheaper generator.

\paragraph{Disadvantages.} The whole system is capped by retrieval quality --- if the right
passage is not retrieved, the answer will be wrong; added latency and infrastructure;
chunking strategy matters a great deal; and the model may still ignore or contradict the
retrieved context.

\paragraph{RAG versus fine-tuning.}
\begin{center}
\begin{tabularx}{\textwidth}{lYY}
\textbf{Dimension} & \textbf{RAG} & \textbf{Fine-tuning} \\
\hline
Knowledge freshness & Immediate --- update the index & Stale --- requires retraining \\
Transparency & High --- the retrieved sources can be shown & Low --- knowledge is diffused
into weights \\
Compute cost & Low up-front, higher per query (longer prompts) & High up-front, cheaper per
query \\
Hallucination risk & Lower, because answers are grounded in retrieved text & Unchanged or
worse; fine-tuning on facts can teach confident invention \\
Best at teaching & \emph{Knowledge} --- what to say & \emph{Behaviour} --- how to say it
(style, format, tone) \\
\end{tabularx}
\end{center}

\paragraph{The classwork question, answered.} \emph{You need a model that imitates the writing
style of a particular newspaper while also being factual --- which solution?} Both, and the
split follows the last row of the table. Fine-tune (or LoRA) on that newspaper's archive to
acquire its \emph{style}; use RAG to supply the \emph{facts} for each specific article. Style
is behaviour, so it belongs in the weights; facts are knowledge, so they belong in the context.
Fine-tuning alone would produce a model that writes beautifully in house style and invents its
facts.

\subsection{Reasoning Models and Agentic AI}

\paragraph{Standard LLMs are ``System 1'' thinkers.} Fast, automatic, pattern-matching; one
forward pass yields one answer; autoregressive, with no ability to backtrack, revise or plan.
Consequently they often fail at multi-step mathematics, logic puzzles, planning and
self-correction.

\paragraph{Reasoning LLMs.} A reasoning model is an LLM fine-tuned to break complex problems
into smaller steps --- often called \emph{reasoning traces} --- before producing a final
output. The distinction to hold onto is \emph{pattern matching versus deliberate reasoning}.
CoT prompting asks for this behaviour at inference time; a reasoning model has been trained to
do it natively, typically with reinforcement learning against verifiable answers. Reasoning
LLMs play a critical role in agentic AI.

\paragraph{The agent loop.} At each iteration the agent assembles context from the available
inputs, invokes an LLM to reason and select an action, executes that action, observes the
outcome, and feeds the observation back into the next iteration. The loop repeats until the
task is complete or a stopping condition is reached.

\paragraph{Why the loop is the key idea.} It converts a one-shot text generator into a system
that can act, observe consequences and correct itself. The model gains a feedback signal from
the world, which is exactly what a single forward pass lacks.

\paragraph{Convergence.} Every major AI organisation has converged on the same execution
pattern, summarised by Lilian Weng as
\[
\text{Agent} = \text{LLM} + \text{Memory} + \text{Planning} + \text{Tool Use}.
\]
Memory supplies persistence beyond the context window; planning decomposes the goal; tool use
provides grounding (search, code execution, databases, APIs). The Model Context Protocol (MCP)
standardises how tools are exposed to models, so that any tool can be plugged into any agent
without bespoke integration.

\paragraph{Disadvantages.} Errors compound across iterations; loops can run away in cost;
tool-calling introduces security concerns such as prompt injection; and evaluation is very
hard, because success is defined over a whole trajectory rather than a single output.

\section{Part X: Evaluation and Responsibility}
\label{sec:eval}

\subsection{Why Evaluating LLMs Is Hard}

\begin{itemize}
\item Generative models produce open-ended responses, so correctness is often \textbf{subjective}
  --- there is no single reference answer to compare against.
\item Alignment is subjective as well: how helpful a model is depends on the user.
\item \textbf{Benchmark contamination:} massively trained LLMs may already have been trained
  on the very data used to test them, so a high score may measure memorisation rather than
  ability.
\end{itemize}
Add to these that capabilities are unevenly distributed (a model can pass a bar exam and fail
at counting letters), and that benchmark scores are sensitive to prompt formatting.

\subsection{Common Benchmarks}

\begin{center}
\begin{tabularx}{\textwidth}{lY}
\textbf{Benchmark} & \textbf{What it measures} \\
\hline
MMLU & Multitask accuracy across 57 subject areas \\
HumanEval & Code generation \\
BIG-Bench & Diverse reasoning tasks \\
GSM8K & Grade-school mathematics word problems \\
MT-Bench & Multi-turn conversation and reasoning \\
\end{tabularx}
\end{center}
Benchmarks are useful but not sufficient. They are cheap, comparable and reproducible; they
are also gameable, contaminated over time, and silent about the qualities users actually care
about.

\subsection{LLM-as-a-Judge}

Extremely popular at present: if prompted properly, a strong LLM scoring another model's
output is a scalable, flexible and surprisingly accurate alternative to human evaluation.

\begin{center}
\begin{tabularx}{\textwidth}{YY}
\textbf{Advantages} & \textbf{Disadvantages} \\
\hline
Fast and scalable & Verbosity bias --- longer answers are scored higher \\
Aligns reasonably well with human preferences & Egocentric bias --- a judge prefers text
resembling its own output \\
Flexible --- any criterion can be written into the rubric & Cost --- an extra model call per
evaluation \\
Consistent --- no annotator fatigue or drift & Ordering bias --- position of A and B affects
the verdict \\
Explainable --- the judge can state its reasons & Trustworthiness --- the stated reasons may
not be the real ones \\
\end{tabularx}
\end{center}
Practical mitigations: swap the order of candidates and average; use a fixed rubric; and
calibrate the judge against a human-scored sample.

\subsection{Human Evaluation}

The gold standard, but expensive. For some qualities --- humour, empathy, tone --- it may be
the only valid method. It is also inconsistent and subjective across annotators, and scaling
it is a nightmare. In practice the field uses a pyramid: automatic benchmarks for fast
iteration, LLM-as-judge for breadth, human evaluation for final decisions.

\subsection{Responsible LLMs}

\begin{itemize}
\item \textbf{Hallucinations.} Fluency is optimised, truth is not. Mitigations: RAG,
  citations, escape hatches, calibration, human review for high-stakes use.
\item \textbf{Bias and fairness.} Models reproduce and can amplify the distribution of their
  training data --- gender, race, religion, language and geography. Mitigations: data
  auditing, targeted evaluation, alignment training. Note that debiasing is itself a value
  judgement about which distribution is correct.
\item \textbf{Privacy and intellectual property.} Training data may contain personal
  information and copyrighted work; models can memorise and regurgitate rare strings; and the
  legal status of training on scraped text remains contested.
\item \textbf{Safety and security.} Misuse for disinformation, fraud and malware; jailbreaks
  and prompt injection; over-reliance by users on confident wrong answers; and, in agentic
  settings, real actions taken in the world.
\end{itemize}

\subsection{Future Directions}

Long-horizon models that maintain coherence over extended tasks; personalised assistants;
on-device LLMs (which is what Part VII's efficiency methods are ultimately for); stronger
reasoning; autonomous agents; and alignment and interpretability, so that we can verify what
these systems are actually doing.

\section{Appendix A: Points Where the Slides Need Care}
\label{sec:pitfalls}

\begin{enumerate}
\item \textbf{Grouped-query attention is listed under ``LLaMA 1 new concepts''.} In the
  published papers, LLaMA 1 used standard multi-head attention. GQA appeared in
  \textbf{LLaMA 2}, and only in the 34B and 70B variants. If a question asks which LLaMA
  generation introduced GQA, the correct answer is LLaMA 2.
\item \textbf{RMSNorm appears in the LLaMA architecture diagram but not in the bulleted list
  of new concepts,} where the item is simply ``Normalization''. RMSNorm is one of LLaMA's four
  defining changes and should be named explicitly alongside RoPE, GQA and SwiGLU.
\item \textbf{GPT-1's parameter count} is given as 120 million in the slides; the paper's
  configuration is usually reported as about 117 million. Either figure is acceptable; do not
  build an argument on the difference.
\end{enumerate}

\section{Appendix B: One-Page Comparison Tables}

\subsection{Recurrent versus Attention-Based}

\begin{center}
\begin{tabularx}{\textwidth}{lYYY}
\textbf{Property} & \textbf{Simple RNN} & \textbf{LSTM / GRU} & \textbf{Transformer} \\
\hline
Path length between distant tokens & $O(n)$ & $O(n)$ & $O(1)$ \\
Parallel over positions & No & No & Yes \\
Compute per layer & $O(n \cdot d^2)$ & $O(n \cdot d^2)$ & $O(n^2 d + n d^2)$ \\
Long-range dependencies & Poor & Moderate & Strong \\
Memory of past & One vector & Cell state plus hidden state & Full context, re-read each step \\
Inductive bias & Strong (recency) & Strong (recency, gated) & Weak --- needs much more data \\
Best regime & Tiny data, streaming & Small--medium data & Large data plus GPUs \\
\end{tabularx}
\end{center}

\subsection{Which Efficiency Method for Which Constraint}

\begin{center}
\begin{tabularx}{\textwidth}{lY}
\textbf{Your constraint} & \textbf{Reach for} \\
\hline
Model will not fit in VRAM & Quantization (and, if available, a smaller distilled model) \\
Cannot afford to fine-tune all parameters & PEFT / LoRA \\
Need many task variants served from one base & LoRA adapters swapped at serving time \\
Need a genuinely smaller and faster model & Distillation, or structured pruning \\
Need more knowledge capacity at fixed FLOPs & Mixture of Experts \\
Need current or private facts & RAG (not fine-tuning) \\
Need a particular style or format & Fine-tuning (not RAG) \\
\end{tabularx}
\end{center}

\section{Appendix C: Rapid Revision Questions}

\begin{enumerate}
\item \textbf{Why does the LSTM solve the vanishing gradient problem?} Because the cell state
  update is additive, $C_t = f_t \odot C_{t-1} + i_t \odot \tilde C_t$, so the gradient through
  time is multiplied by the forget gate rather than by a chain of $\tanh$ derivatives. With
  $f_t \approx 1$ the gradient passes back nearly unchanged.

\item \textbf{Why divide by $\sqrt{d_k}$ in attention?} Dot products of high-dimensional
  vectors have large variance; large logits saturate the softmax, whose gradient then vanishes.
  Scaling keeps logits in a range where softmax remains trainable.

\item \textbf{Why does a Transformer need positional encoding when an RNN does not?}
  Self-attention is permutation-invariant --- it treats the input as a set. An RNN encodes
  order implicitly through the sequential order of its updates.

\item \textbf{Why can BERT not generate text?} It is an encoder trained with bidirectional
  attention. Any position can see every other, so there is no way to define a next-token
  prediction that does not leak the answer, and there is no autoregressive decoder.

\item \textbf{Why does BERT use 80/10/10 masking?} \texttt{[MASK]} never occurs at fine-tuning
  time, so training on 100\% \texttt{[MASK]} creates a pretrain/fine-tune mismatch. Random and
  unchanged tokens force good representations of ordinary tokens too, which matters most in the
  frozen feature-extraction setting.

\item \textbf{What is the single most important claim of the Chinchilla paper?} For a fixed
  compute budget, model size and training data should scale equally; most prior models were
  too large for their data, and doubling data helps at least as much as doubling parameters.

\item \textbf{Why did GQA appear?} GPUs compute much faster than they move memory, so
  autoregressive inference is bandwidth-bound, and the KV cache dominates that bandwidth.
  Sharing keys and values across groups of query heads shrinks the cache with little quality
  loss.

\item \textbf{Why is RoPE preferred over sinusoidal encoding?} It encodes relative position
  through rotation, so attention scores depend only on the offset between positions; it is
  applied inside every layer rather than once at the input; it adds no parameters; and it can
  be extended to longer contexts by rescaling frequencies.

\item \textbf{What does a MoE model buy you?} Knowledge capacity decoupled from per-token
  compute: many parameters exist, few are active. The cost is that all of them must still be
  held in memory, plus routing instability.

\item \textbf{Difference between quantization, pruning, PEFT and distillation?} Quantization
  reduces the precision of weights; pruning removes weights or structures; PEFT reduces the
  number of \emph{trainable} weights without shrinking the model; distillation trains a new,
  smaller model to imitate a large one.

\item \textbf{Why RLHF rather than more supervised fine-tuning?} SFT can only imitate
  demonstrations, so it teaches format, not preference, and cannot express that one valid answer
  is better than another. Preference comparisons are easier for humans to give and support
  optimisation beyond what any demonstrator wrote.

\item \textbf{RAG or fine-tuning?} RAG for knowledge that is fresh, private, large or must be
  cited; fine-tuning for behaviour --- style, format, tone, task-specific output structure.
  Combine them when you need both.

\item \textbf{Why does chain-of-thought help?} It gives the model more computation for a hard
  problem (each generated token is another forward pass, and the text acts as working memory)
  and conditions it on the kind of context that, in the training distribution, precedes correct
  answers.

\item \textbf{What is the biggest weakness of long-context models?} They do not use the whole
  window equally --- the ``lost in the middle'' effect --- and attention cost grows
  quadratically with length.

\item \textbf{Name the four architectural differences between LLaMA and the original
  Transformer decoder.} Pre-norm RMSNorm instead of post-norm LayerNorm; rotary positional
  encoding instead of sinusoidal; grouped-query attention with a KV cache instead of plain
  multi-head attention (LLaMA 2, larger variants); and a SwiGLU gated feed-forward network
  instead of a ReLU MLP.
\end{enumerate}

\end{document}